\documentclass[12pt, a4paper]{article}

\usepackage{amsmath}
\usepackage{graphicx}
\usepackage{booktabs}
\usepackage{hyperref}

\title{Ecological Footprint and Degrowth}
\author{W}
\date{\today}

\begin{document}

\maketitle

\tableofcontents

\section{Introduction}
Degrowth challenges the assumption that GDP growth is necessary for human well-being \cite{kallis2018}. The relationship between economic output $Y$ and ecological impact $E$ can be expressed as:

\begin{equation}
    E = P \times A \times T
\end{equation}

where $P$ is population, $A$ is affluence, and $T$ is technology intensity.

\section{Key Indicators}

\subsection{Planetary Boundaries}
The following table summarizes selected planetary boundaries:

\begin{table}[h]
\centering
\begin{tabular}{lcc}
\toprule
\textbf{System} & \textbf{Boundary} & \textbf{Current Status} \\
\midrule
CO\textsubscript{2} concentration & 350 ppm & Exceeded \\
Biodiversity loss & 10 E/MSY & Exceeded \\
Nitrogen cycle & 35 Tg N/yr & Exceeded \\
\bottomrule
\end{tabular}
\caption{Selected planetary boundaries}
\label{tab:boundaries}
\end{table}

\subsection{The Degrowth Equation}


\section{Conclusion}
As shown in Table \ref{tab:boundaries}, multiple boundaries have already been transgressed. Degrowth offers a framework for operating within these limits.

\begin{thebibliography}{9}
\bibitem{kallis2018}
Kallis, G. (2018). \textit{Degrowth}. Agenda Publishing.
\end{thebibliography}

\end{document}